% concatenated from rings.tex
\documentclass{amsart}
%\documentclass{birkau}
\usepackage[parfill]{parskip} 
\usepackage{amsmath,amssymb,amsthm}
\usepackage[pdftex]{hyperref}
%\usepackage[margin=4cm]{geometry}
\usepackage{euscript}
\usepackage{tikz-cd}
\usepackage{appendix}
\usepackage{comment}
\usepackage{mathtools}
%\usepackage{mathabx} %%to do updownarrows 
\usepackage{graphicx} %% to do parallelogramm
%\usepackage[scr=boondox]{mathalpha} %%to do lowercase calligraphic 
\usepackage{booktabs}

\hypersetup{
  colorlinks,
  citecolor=black,
  filecolor=black,
  linkcolor=black,
  urlcolor=black
}

\numberwithin{equation}{section}

\theoremstyle{plain}
\newtheorem{theorem}{Theorem}[section]
\newtheorem*{theorem*}{Theorem}
\newtheorem{lemma}[theorem]{Lemma}
\newtheorem{proposition}[theorem]{Proposition}
\newtheorem{corollary}[theorem]{Corollary}
\newtheorem{example}[theorem]{Example}


\theoremstyle{definition}
\newtheorem{definition}[theorem]{Definition}
\newtheorem{notation}[theorem]{Notation}
\newtheorem{remark}[theorem]{Remark}

\DeclareMathOperator{\Eq}{Equiv}
\DeclareMathOperator{\Stone}{Stone}
\DeclareMathOperator{\op}{op}
\DeclareMathOperator{\Mat}{Mat}
\DeclareMathOperator{\Alg}{Alg}
\DeclareMathOperator{\id}{id}
\DeclareMathOperator{\pr}{pr}
\DeclareMathOperator{\supp}{supp}
\DeclareMathOperator{\fin}{fin}


%
\begin{document}
%
\title{Rings and Boolean Algebras as Algebraic Theories}

%\titlerunning{Abbreviated paper title}
% If the paper title is too long for the running head, you can set
% an abbreviated paper title here
%
\author{Arturo De Faveri} %\inst{1}
%\orcidID{0000-1111-2222-3333}}
%
%\authorrunning{Arturo De Faveri}
% First names are abbreviated in the running head.
% If there are more than two authors, 'et al.' is used.
%
\email{defaveri@irif.fr}
\address{IRIF\\
CNRS and Universit\'e Paris Cit\'e\\
8 Place Aur\'elie Nemours, 75013 Paris, France
}

\begin{abstract}
%The abstract should briefly summarize the contents of the paper in
%150--250 words.
We present a unified framework for representing commutative rings through affine algebraic theories and 
Boolean rings through hyperaffine algebraic theories.
This yields categorical equivalences between these theories and, respectively, commutative rings and Boolean rings. 
%a new analysis of certain classes of modules over Boolean rings.
We then analyse models of affine theories over a Boolean ring $B$, comparing them with the models of hyperaffine theories, the well-known $B$-sets. 
Two novel characterisations are presented: the first defines these models as Boolean vector spaces equipped with an action of the Boolean ring; 
the second provides a representation in terms of sheaves, in analogy with $B$-sets.
Finally, we establish a connection between hyperaffine theories and multidimensional Boolean algebras, a recently introduced generalisation of Boolean algebras.   
%The resulting structures naturally capture the algebraic semantics of the if-then-else construct in programming languages.
%\keywords{Algebraic theories \and Commutative rings \and Boolean rings \and Affine spaces \and Boolean actions}

\end{abstract}

\maketitle



\section{Introduction}
An \emph{algebraic theory} (also called a \emph{clone}) is an abstraction of the notion of algebraic operations and the equations they satisfy.
In universal algebra, an equational class of algebras is specified by a signature of operations and a set of equational identities.
An algebraic theory $T$ encodes the same information in different terms: it is a collection of abstract operations of finite arities together with a rule for their composition.  
%for each $n$, it specifies all possible $n$-ary operations, and composition allows building new operations from existing ones 
% it is a category $\mathbb{T}$ whose objects can be thought of as finite powers of a single “generic” object, and whose morphisms represent abstract operations of various arities.
An \emph{algebra} of the theory $T$ is then a set equipped with actual operations corresponding to the elements of the theory, in such a way that composition is respected.
%In this sense, algebraic theories provide a framework for describing all classical algebraic structures without fixing a particular underlying set in advance.

Given an algebraic structure, a natural and intriguing question is how to represent it as an algebraic theory.
In this paper we focus on rings, and we ask whether there exists a full embedding $F: \mathbf{Ring} \to \mathbf{AlgTh}$ of the category of rings into that of algebraic theories such that suitable abstract properties of the theory $F(R)$ enable to reconstruct perfectly the ring $R$. 
% -> aggiungere: voglio anche tornare indietro dalle proprietà delle teoria 
The first idea that comes to mind is to associate with $R$ the theory $F(R)$ of (left) $R$-modules. 
The abstract operations, in this case, are the tuples $(r_1,\ldots,r_n) \in R^n$, thought of as linear combinations $r_1x_1+\cdots+r_nx_n$. 
If the ring is commutative, then all the abstract operations are commutative.
%according to Definition \ref{def:icd}.  
Two binary operations $s,r$ are said to commute if 
\begin{equation}
  \label{eq:comm}
r(s(x,y),s(z,w)) = s(r(x,z),r(y,w)) \text{,}
\end{equation}
and this condition can be generalised to higher arities. 
%Despite this being a full embedding, 
However, if the goal is to reconstruct $R$ from abstract properties like \eqref{eq:comm}, the encoding of $R$ via $R$-modules is not sufficiently rich. 

It turns out that there are two orthogonal constructions that allow us to reconstruct $R$ perfectly
and these two constructions have in common the idea of encoding the ring through the binary operations of the corresponding theory. 
The first, which works in the case of a commutative ring, associates with $R$ the theory of affine $R$-modules, namely the theory whose operations are linear combinations $r_1x_1+\cdots+r_nx_n$ such that $r_1+\cdots+r_n=1$. 
The operations of this theory are commutative and idempotent.
% which in the binary case takes the form
%\begin{equation}
%\label{eq:id}
%r(x,x) = x \text{.}
%\end{equation}
%The intuition consists in imagining an abstract binary operation as an affine combination $rx+(1-r)y$. 
The second construction, specific to Boolean rings, associates with a Boolean ring $B$ the theory of ``hyperaffine'' $B$-modules, that is, the theory whose operations are linear combinations $b_1x_1+\cdots+b_nx_n$ such that $b_1+\cdots+b_n=1$ and $b_i b_j =0$ for $i \neq j$. 
In this case, a strong property holds
\begin{equation}
%\label{eq:split}
b(b(x,y), b(z,w)) = b(x,w) \text{,}
\end{equation}
making each binary operation a rectangular band. 
These observations can be crystallised in the following statement:
for every commutative (resp. Boolean) ring $R$, the theory of affine (resp. ``hyperaffine'') $R$-modules is affine (resp. hyperaffine); 
moreover, any affine (resp. hyperaffine) theory is uniquely determined by the ring of its binary operations 
%(cf. Theorems 3.7 and 3.14).
(cf. Theorems \ref{thm:equibool} and \ref{thm:equiring}).

\begin{comment}
One of the reasons that make this approach interesting is that studying the algebras of affine and hyperaffine theories can provide an elegant aetiology for the axiomatisation of if-then-else construct in programming languages.
For example, among the axioms for if-then-else originally introduced by McCarthy \cite{McC63}, one corresponds to idempotence (\eqref{eq:id} above), others encode the laws of composition and unitality in a clone, others express \eqref{eq:split}, and another one is commutativity \eqref{eq:comm}.
Given a Boolean ring $B$, a set together with an action of $B$ satisfying these conditions is what Bergman called a $B$-set \cite{B91}. 
%It is natural to ask how these axioms could be weakened to capture alternative or more nuanced behaviors of the if-then-else operator.
\end{comment}
%In questo articolo ci proponiamo di sviluppare le due costruzioni sotto un cappello comune. 

%\subsubsection{Contributions.}
Our first main contribution is the development of a convenient framework in which the ``hyperaffine'' and ``affine'' constructions can be realised.
%It is within this framework that we obtain Theorems~\ref{thm:equibool} and~\ref{thm:equiring}.
We first 
%provide a streamlined account of results in \cite{J90,G24}, 
provide a new proof that hyperaffine theories form a category equivalent to that of Boolean rings.
This result was originally stated in \cite[Theorem 4.9]{G24}.
We then apply the same strategy to show in Theorem~\ref{thm:equiring} that affine theories correspond precisely to commutative rings.
We highlight the importance of the role played by the \emph{coefficients} of a theory to the end of presenting both constructions in a single, coherent setting.
%Although some ingredients already appear in the literature, to the best of our knowledge there is no source where is stated in the form adopted here;
%nor, more importantly, where both constructions are presented within a single, coherent way.
%We believe that this unified perspective clarifies the parallels between the two settings and opens the way to further generalisations.
%While the former is specific to Boolean rings, the other applies to all commutative rings.
%Although Theorem \ref{thm:equibool} appears in \cite{G24}, to the best of our knowledge, despite all the necessary ingredients being present in the literature, there is no source in which Theorem~\ref{thm:equiring} is expressed in the form adopted in this paper, nor, more importantly, where the two constructions are presented in a single, uniform framework.
%We hope that bringing these insights together in one place will be useful.
Given a Boolean ring $B$, as $B$ is commutative, we have two ways of encoding $B$ as an algebraic theory: either as a hyperaffine theory or as an affine theory. 
Models of hyperaffine theories are known to be $B$-sets \cite{B91}, so that we are led to the question: what are the models of an affine theory whose ring of coefficients is Boolean? 
%A_B. Models of H_B are the well-known B-sets. 
%We therefore investigate models of an affine theory when the coefficient ring is Boolean.
% a task we undertake in this section.
Our second main contribution is the answer to this question.
%the study of the models of an affine theory over a Boolean ring $B$.
% in a manner analogous to the investigation of hyperaffine theories.
We provide two characterisations for the models of an affine theory over a Boolean ring $B$.
The first, Theorem \ref{thm:new1}, is purely algebraic:
%and parallels the definition of a $B$-set: 
the models of an affine theory over a Boolean ring are Boolean vector spaces equipped with a compatible action of $B$ which is idempotent and commutative. 
%One axiom defining the action--namely, the one incompatible with the affine structure--is dropped, but the set is endowed with a compatible structure of a Boolean vector space.
The second, Theorem \ref{thm:new2}, shows that these structures admit a sheaf representation analogous to $B$-sets (see \cite{B91}).
%If $B$-sets provide the standard semantics for the if-then-else operator, we speculate that affine $B$-modules could be applied in a semantics of if-then-else in linear contexts, although a thorough investigation of this is left for future work.


Inspired by the representation of Boolean algebras through a ternary ``conditional disjunction" operation (if-then-else), some authors have recently introduced $n$-dimensional Boolean algebras ($n$BAs), featuring an $(n+1)$-ary operation $q$ and $n$ constants $e_1, \ldots, e_n$ \cite{BS,BLPS18,SBLP20}. 
These structures provide the algebraic semantics for a many-valued logic with symmetric truth values $e_1, \ldots, e_n$. 
We highlight a connection between hyperaffine algebraic theories and $n$BAs. 
We prove in Theorem \ref{thm:nba} that every hyperaffine theory $T$ naturally induces an $n$BA structure on the set of its $n$-ary operations; 
conversely, a coherent sequence of $n$BAs uniquely determines a hyperaffine theory. 
This link originates from the dual interpretation of the $(n+1)$-ary operator $q$, viewed on the one hand as a composition operator, and on the other as a generalised if-then-else.
In \cite{BS}, it is shown that $n$BAs are related to skew Boolean algebras \cite{Leech1}.  
These considerations suggest that, in the affine case, the corresponding theories naturally point toward the introduction of suitable notions of $n$-dimensional commutative rings and skew rings, in analogy with the Boolean case.

\subsection{Related works.}
Hyperaffine theories were introduced by Johnstone \cite{J90} in the context of a theorem providing necessary and sufficient syntactic conditions for a variety of algebras to form a cartesian closed category.
This result has since been elegantly revisited, opening up new perspectives and developments \cite{G24,G25}.
Among other contributions, \cite{G24} offers a detailed account of the connection between hyperaffine theories and their models, namely $B$-sets.
More about $B$-sets can be found in \cite{S98a,S98b}.
On the other hand, models of affine theories essentially capture what the Polish school calls Mal'cev modes \cite{RS02},
culminating a long line of research on varieties of affine modules \cite{SO66,C75,S77,PRS95}.
%In turn, reducts of models of these theories are algebras known as convex spaces or barycentric algebras \cite{S49,N70,KRS19},
%which play a significant role in algebraic accounts of probability \cite{K08,J10}.
The models of hyperaffine theories are connected with the algebraic treatment of the if-then-else construct, a topic that has been extensively studied in the semantics of programming languages \cite{M90,M92,M93,JS09}. 
%\cite{GM87,MN87,M90,M92}.
%The clearest account of the semantic behaviour of if-then-else is arguably Bergman's formulation \cite{B91}, based on the notion of an action of a Boolean ring on a set.
%Alternative approaches have been proposed in which the Boolean ring $B$ acts not on sets but on other algebraic structures: join-semilattices with bottom \cite{M93}, where programs are viewed as ordered structures, or semigroups \cite{JS09}, where programs are viewed as functions.
% thus creating a bridge with Kleene algebras with tests \cite{K97}.

\subsection{Outline of the paper.}
Section~\ref{sec:alg} introduces the preliminary material on algebraic theories.
Section~\ref{sec:repr} is divided into two main parts.
In the first, it is shown that hyperaffine theories form a category equivalent to that of Boolean rings.
In the second part, we apply the previously developed blueprint to show that affine theories, by contrast, correspond to commutative rings.
Section \ref{sec:nba} explores the link between hyperaffine theories and $n$-dimensional Boolean algebras.
% algebraic structures giving the semantics of a symmetric many-valued logic. 
Section~\ref{sec:spaces} is devoted to the study of the models of these theories, with particular attention to the models of affine theories when the ring of coefficients is Boolean.
% and to a comparison with the models of hyperaffine theories.
%Finally, Section~\ref{sec:conclusions} concludes the paper with a summary of the main results and some directions for future research.
Finally, Section \ref{sec:conclusions} wraps up with concluding remarks and directions for future work.

\section{Varieties and algebraic theories}
\label{sec:alg}
In a broad sense, a variety is the category of models of a finitary and one-sorted equational theory. 
%For theories with a mere set of function symbols, free V-structures always exist; we will relax this by allowing a proper class of function symbols, but still assuming that free V-structures exist. So the category of complete join-lattices is a variety in our sense, but not the category of complete Boolean algebras. We write Var for the category of varieties and concrete functors between them.
%A variety is non-degenerate if it contains a model with at least two elements. 
%Up to isomorphism, there are two degenerate varieties:  is the full subcategory of Set on the one-element sets, while V2 is the full subcategory on the zero- and one-element sets. The former is the category of models of the equational theory with no operation symbols and the axiom x = y; while the latter is the category of models of the theory with a single constant c and the axiom c = x.
A given variety may be axiomatised by operations and equations in many ways; 
however, there is always a canonical choice, which is captured by the notion of \emph{algebraic theory} due to Lawvere \cite{Law63} (see also \cite[Section 3]{B2} and \cite[Chapter 11]{ARV10}). 
Here we give a formulation of algebraic theory essentially equivalent to what universal algebraists call an abstract clone \cite{T93}.
% and what a category theorist would call a monad relative to the identity functor Set → Set.

\begin{definition}
  An \emph{algebraic theory} $T$ is a family of sets $T(n)$, for each $n \in \mathbb{N}$, together with elements $\pi^n_i \in T(n)$, $1 \le i \le n$, and 
a composition operator $ \circ : T(n) \times T(m)^n \to T(m)$, for each $m,n \in \mathbb{N}$. 
These data have to satisfy, writing $f(g_1, \ldots, g_n)$ in place of $f \circ (g_1, \ldots, g_n)$, two unit laws for $1 \le i \le n$
\begin{equation}
  \label{eq:c12}
  \pi^n_i(g_1, \ldots, g_n) = g_i \quad \text{ and } \quad  f(\pi^n_1, \ldots, \pi^n_n) = f 
\end{equation}
and the associative law 
\begin{equation}
  \label{eq:c3}
  f(g_1(h_1,\ldots h_m),\ldots,g_n(h_1,\ldots,h_m)) = f(g_1,\ldots,g_n)\circ (h_1,\ldots h_m)\text{.}
\end{equation}
\end{definition}

We call an element of $T(n)$ an $n$-ary \emph{operation} in $T$. 

\begin{definition}
  A morphism of theories $\varphi: T \to T'$ comprises, for every $n$, a function $\varphi: T(n) \to T'(n)$ satisfying
\begin{equation}
  \varphi(\pi^n_i) = \pi^n_i \text{ for } 1 \le i \le n  \text{ and } \varphi(f(g_1, \ldots, g_n)) = \varphi(f)(\varphi(g_1), \ldots, \varphi(g_n))
\end{equation}
for all $f \in T(n)$ and $g_1, \ldots, g_n \in T(m)$. 
\end{definition}

%Given a category $\EuScript{C}$ with finite products we can define a theory $T_{\EuScript{C},x}$ by letting $T_{\EuScript{C},x}(n):=\EuScript{C}(x^n,x)$, and it can be shown that every algebraic theory is actually of this form. 


\begin{definition}
A \emph{model} or an \emph{algebra} for an algebraic theory $T$ is a set $X$ 
together with, for each $n \in \mathbb{N}$, a left action $\cdot : T(n) \times X^n \to X$ satisfying two conditions: 
%if $\pi^n_i$ are the $n$ projections, then   
  \begin{equation}
    \label{eq:alpha1}
    \pi^n_i \cdot (x_1,\ldots,x_n) = x_i 
  \end{equation}
for all $1 \le i \le n$, and 
  \begin{equation}
    \label{eq:alpha2}
    f(g_1, \ldots, g_n) \cdot (x_1, \ldots, x_m) = f \cdot (g_1 \cdot (x_1, \ldots, x_m), \ldots, g_n \cdot (x_1, \ldots, x_m))
  \end{equation}
  for each $f \in T(n), g_1, \ldots, g_n \in T(m), x_1, \ldots, x_m \in X$.
  The category of models of a theory is called an \emph{algebraic variety}. 
\end{definition}

For each $m$, the set $T(n)$ is a model with action 
\begin{equation*}
  \cdot: T(n) \times T(m)^n \to T(n) \quad \text{ given by } \quad \circ : T(n) \times T(m)^n \to T(n)\text{.} 
\end{equation*}
We call $T(m)$ the free model on $n$ generators.
%Given an algebraic theory $T$, the free model on $n$ generators, $n \in \mathbb{N}$, is given by $T(n)$ itself. 
There is an obvious forgetful functor from the category of models of $T$ to the category of sets, associating with the model $X$ its underlying set $X$.
% giving the \emph{underlying set} of $X$.  
The underlying set of the free model on $n$ generators is $T(n)$. 
In this light, we can freely identify the operation $f \in T(n)$ with the element $f(x_1, \ldots, x_n)$ in the free model on the $n$ generators $x_1, \ldots, x_n$.


The initial theory $S$ is such that $S(n)$ contains only the $n$ projections $\pi^n_i$ for $n \ge 1$ and $S(0)$ is empty; models of $S$ are mere sets. 
The terminal theory $U$ is given by $U(n)=\{*\}$ for all $n$ and its only model is the singleton. 
The terminal theory has a subtheory $U'$ (i.e., there is an injective morphism $U' \to U$) given by 
\begin{equation*}
    U'(n) = 
    \begin{cases}
      \varnothing & \text{ if } n = 0 \\
      \{*\} & \text{ if } n \neq 0 \text{,}
    \end{cases}
  \end{equation*}
whose models are either the empty set or the singleton. 
We refer to $U$ and to $U'$ as the two \emph{degenerate} theories. 

\subsection{Mal'cev operations, affine and hyperaffine theories.}
We recall the following well-known definitions. 
%The first one dates back to \cite{Law68}. 

\begin{definition} 
  \label{def:icd}
  Let $T$ be an algebraic theory; $f \in T(n)$ is said to 
  \begin{itemize}
  \item be \emph{idempotent} if $f(x,\ldots,x)=x$; 
  \item \emph{split} if 
  %\begin{equation*}
    $f(f(x^1_1, \ldots x^1_n), \ldots, f(x^n_1, \ldots, x^n_n)) = f(x^1_1, \ldots, x^n_n)\text{.}$
  %\end{equation*}
  \end{itemize}
  An idempotent, splitting operation is called a \emph{decomposition operation}. 
  Moreover, $f \in T(n)$ and $g \in T(m)$
  are said to \emph{commute} if 
    \begin{equation*}
      f(g(x^1_1, \ldots x^1_m), \ldots, g(x^n_1, \ldots x^n_m)) = g(f(x^1_1, \ldots x^n_1), \ldots, f(x^1_m, \ldots x^n_m))\text{.}
    \end{equation*}
  A theory is \emph{idempotent} (resp. \emph{commutative}) if every operation is (resp. if every pair of operations commute). 
\end{definition}

%An algebraic variety whose theory is idempotent and commutative is called a variety of \emph{modes} \cite{RS02}. 
Another essential ingredient in the paper is that of a Mal'cev operation. 
Roughly speaking, a Mal'cev operation in a commutative theory generalises the operation $f(x,y,z)=x-y+z$ in Abelian groups. 
The importance of this concept lies in the fact that the presence of a Mal'cev operation in a theory is reflected in the internal structure of its models: 
given two congruences $R$, $S$ on a model $X$ of $T$, their composition $R \circ S$ commutes, meaning that $R \circ S = S \circ R$, iff $T(3)$ has a Mal'cev operation.
% (see \cite[Theorem 2.2.2]{BB04}). 
Formally, a Mal'cev operation is defined as follows; 
for an exhaustive and much more general treatment we refer the reader to \cite{Smith,BB04}. 

\begin{definition} 
  Let $T$ be an algebraic theory. 
  An element $p \in T(3)$ is called a \emph{Mal'cev operation} if 
  \begin{equation*}
    x=p(x,y,y) \quad \text{ and } \quad p(x,x,y)=y \text{.}
  \end{equation*}
\end{definition}

%The following lemma is straightforward. 

\begin{lemma}
  \label{lem:asscomm}
  Let $p \in T(3)$ be a Mal'cev operation that commutes with itself. 
  Then $p$ satisfies 
    \begin{description}
    \item[M1] $p(x,y,p(t,u,v)) = p(p(x,y,t),u,v)$; 
    \item[M2] $p(x,y,z) = p(z,y,x)$. 
  \end{description}
  %Moreover, \emph{(M1)} holds iff
  %\begin{equation*}
  %  p(x,y,p(y,u,v)) =p(x,u,v) \quad \text{ and } \quad p(p(x,y,u),u,v) =p(x,y,v) \text{.}
  %\end{equation*}
  Moreover, if $p$ and $p'$ are two Mal'cev operations in $T(3)$ which commute with each other, then they coincide.  
\end{lemma}
\begin{proof}
  Firstly, we prove (M1)  
  \begin{align*}
    p(x,y,p(t,u,v)) & = p(p(x,z,z),p(y,z,z),p(t,u,v)) \\
    & = p(p(x,y,t),p(z,z,u),p(z,z,v)) \\
    & = p(p(x,y,t),u,v)\text{.}
  \end{align*}
  Then we prove (M2)
  \begin{align*}
    p(x,y,z) & = p(p(z,y,x),p(z,y,x),p(x,y,z)) & \\
    & = p(z,y,p(x,p(z,y,x),p(x,y,z))) & \text{(M1)} \\
    & = p(z,y,p(p(z,z,x),p(z,y,x),p(x,y,z))) & \\
    & = p(z,y,p(p(z,z,x),p(z,y,y),p(x,x,z))) &\\
    & = p(z,y,p(x,z,z))& \\
    & = p(z,y,x)\text{.}& 
  \end{align*} 
  Finally, let $p$ and $p'$ be two Mal'cev operations that commute with each other.
  Then 
  \begin{align*}
    p(x,y,z) & = p(p'(x,y,y),p'(y,y,y),p'(y,y,z)) \\
    & = p'(p(x,y,y),p(y,y,y),p(y,y,z)) \\
    & = p'(x,y,z)
  \end{align*}
  concluding the proof. 
\end{proof}
 
\begin{comment}
The category of algebraic theories is equivalent to the category of finitary monads on sets  \cite{L65}; 
this equivalence commutes with the functors sending a monad to its algebras, and an algebraic theory to its models. 
%For a textbook exposition of this link the reader could consult \cite{BW85}. 
Idempotent theories correspond to (finitary) monads on $\mathbf{Set}$ preserving the terminal object. 
Commutative theories have the desirable property that their categories of models admit the structure of symmetric monoidal closed category, and they correspond to (finitary) commutative monads on $\mathbf{Set}$ as defined by Kock \cite{kock}. 
%The interested reader can find a textbook account of commutative theories in \cite[Section 3.10]{B2}. 
\end{comment}

\begin{definition}
  \label{def:hyper-affine}
  We call \emph{affine}\footnote{Note that this terminology is slightly non-standard and some authors use the word `affine' to denote what we call idempotent theories.} an idempotent, commutative theory $T$ with a Mal'cev operation $p \in T(3)$.  
  Moreover, following \cite[Section 4]{J90} we call \emph{hyperaffine} an idempotent, commutative theory in which every operation splits.
\end{definition}

The two degenerate theories $U$ and $U'$ are affine and hyperaffine.

\begin{remark}
  An algebra whose operations are idempotent and commutative is called a \emph{mode} \cite{RS02}. 
  Therefore, models of an affine theory form an algebraic variety of modes. 
  The same applies to a hyperaffine theory. 
\end{remark}
%In the equivalence between finitary monads and algebraic theories, hyperaffine theories correspond to those finitary monads on $\mathbf{Set}$ that preserve finite products
% and neglecting the degenerate theory $U'$, they correspond to monads preserving finite limits 
%(see \cite[Theorem 2.1]{J97} and subsequent discussion). 

The following distributive property will be useful later. 
%lemma. 
\begin{lemma}
  \label{lem:c5}
  Let $T$ be a hyperaffine theory. 
  Then for $f \in T(k)$ and $g_1, \ldots, g_k \in T(n)$,
  % letting $t_i = f(x^1_i, \ldots, x^k_i)$ for $1 \le i \le n$,  we have
  \begin{equation}
    \label{eq:c5}
    \begin{split}
        f(g_1, \ldots, g_k) & \circ (f(x^1_1, \ldots, x^k_1), \ldots, f(x^1_n, \ldots, x^k_n)) \\
        & =     f(g_1(x^1_1, \ldots, x^1_n), \ldots, g_k(x^k_1, \ldots, x^k_n))\text{.}
    \end{split}
    %f(g_1(t_1, \ldots, t_n), \ldots, g_k(t_1, \ldots, t_n)) = 
  \end{equation}
  Moreover, if an idempotent theory $T$ satisfies \eqref{eq:c5}, then $T$ is hyperaffine. 
\end{lemma}
\begin{proof}
    We prove \eqref{eq:c5} 
    \begin{align*}
        & f(g_1, \ldots, g_k) \circ (f(x^1_1, \ldots, x^k_1), \ldots, f(x^1_n, \ldots, x^k_n)) \\
        & = f(g_1(f(x^1_1, \ldots, x^k_1), \ldots, f(x^1_n, \ldots, x^k_n)), \ldots, g_k(f(x^1_1, \ldots, x^k_1), \ldots, f(x^1_n, \ldots, x^k_n)))\\
        & = f(f(g_1(x^1_1, \ldots, x^1_n), \ldots, g_1(x^k_1, \ldots, x^k_n)), \ldots, f(g_k(x^1_1, \ldots, x^1_n), \ldots, g_k(x^k_1, \ldots, x^k_n)))\\
        & = f(g_1(x^1_1, \ldots, x^1_n), \ldots, g_k(x^k_1, \ldots, x^k_n))\text{.}
    \end{align*}
  Now, assume that $T$ satisfies \eqref{eq:c5}. 
  Firstly, every $f \in T(n)$ splits 
  \begin{align*}
    & f(f(x^1_1, \ldots, x^n_1), \ldots, f(x^1_n, \ldots, x^n_n)) &  \\
    & = f(\pi^n_1, \ldots, \pi^n_n) \circ (f(x^1_1, \ldots, x^n_1), \ldots, f(x^1_n, \ldots, x^n_n)) & \\
    & = f(\pi^n_1(x^1_1, \ldots, x^n_1), \ldots, \pi^n_n(x^1_n, \ldots, x^n_n)) & \eqref{eq:c5} \\
    & = f(x^1_1, \ldots, x^n_n)   \text{.} &
  \end{align*}
  Finally, every $f \in T(n)$ and $g \in T(m)$ commute
  \begin{align*}
    & f(g(x^1_1, \ldots x^1_m), \ldots, g(x^n_1, \ldots x^n_m)) & \\
    & = f(g, \ldots, g) \circ (f(x^1_1, \ldots x^n_1), \ldots, f(x^1_m, \ldots x^n_m)) & \eqref{eq:c5}\\
    & = g(f(x^1_1, \ldots x^n_1), \ldots, f(x^1_m, \ldots x^n_m)) & 
  \end{align*}
  proving the claim. 
\end{proof}


\section{Representing rings as algebraic theories}
\label{sec:repr}
%The convenient umbrella we were referring to in the introduction is that of \emph{matricial} theories \cite{E76}. 
%A \emph{commutative ring} $R$ is an Abelian group $(R,+,0)$ endowed with the structure of a commutative monoid $(R,\cdot,1)$ such that multiplication distributes over addition: $r(a+b)=ra+rb$ for all $a,b,r \in R$.  
%All the rings are assumed to be unital and commutative, and all the ring homomorphisms to preserve the unit.  
A \emph{ring} is a tuple $(R,+,0,\cdot,1)$ such that $(R,+,0)$ is an Abelian group, $(R, \cdot, 1)$ is a multiplicative monoid and multiplication distributes over addition: 
$r(a + b) = ra + rb$ and $(a+b)r=ar+br$ for all $a, b, r \in R$.
%In the following we write ring for ring with unit. 
The ring $R$ is said to be \emph{commutative} if the monoid is, and 
\emph{Boolean} if $r^2 = r$ for every $r \in R$. 
A Boolean ring is necessarily commutative. 

Any ring $R$ gives rise to an algebraic theory $T_R$ as follows. 
For $n \in \mathbb{N}$, let $T_R(n)$ be the set of tuples $(r_1, \ldots, r_n) \in R^n$.
% with $m$ rows and $n$ columns and entries in $R$.
These are the abstract $n$-ary operations of the theory.
% $f \in T(n)$ is a tuple $(r_1, \ldots, r_n) \in R^n$. 
The projections in $T_R(n)$ are given by the standard basis of $R^n$: $(1,0, \ldots, 0), \ldots, (0, \ldots, 0,1)$, 
while composition $\circ : T_R(n) \times T_R(m)^n \to T_R(m)$ is given by matrix multiplication:
\begin{equation*} 
\begin{bmatrix}
      r_1 \\
      \vdots \\
      r_n \\
      \end{bmatrix} \circ  
      \begin{bmatrix}
      a^1_1 & \cdots & a^n_1 \\
      \vdots & \ddots & \vdots \\
      a^1_m & \cdots & a^n_m
      \end{bmatrix} =
      \begin{bmatrix}
      a^1_1r_1 + \cdots + a^n_1 r_n\\
      \vdots \\
      a^1_mr_1 + \cdots + a^n_m r_n 
      \end{bmatrix}
\end{equation*}
%$(A,B) \mapsto B \cdot A$. 
%More about theories, called \emph{matricial}, arising from matrices can be found in \cite{E76}. 
Models of the theory $T_R$ are precisely (left) modules over the ring $R$.
% (an \emph{$R$-module} is an Abelian group $(X,+,0)$ with a compatible left action of $R$: $R \times X \to X$).  
%Modules over the ring $R$ coincides with the algebras of the theory $T_R$. 
%For a proof in the case where $R$ is the commutative ring $\mathbb{Z}$ (and therefore $R$-modules are Abelian groups), see \cite[Example 1.6]{ARV10}. 
\subsection{Hyperaffine theories} 
%We now turn our attention to hyperaffine theories. 
We start by examining the case of hyperaffine theories. 
%However, hyperaffine theories are strong enough to imply that the set of binary operations is not only a ring but a Boolean ring. 
%We recall that a Boolean ring is a ring $B$ whose product is idempotent, meaning that $b \cdot b =b$ for every $b \in B$. 
%We recall also that a Boolean ring is necessarily commutative. 


\begin{definition}
  For a given Boolean ring $B$, we define the algebraic theory $H_B$ whose operations in $H_B(n)$ are $(b_1, \ldots, b_n) \in B^n$ such that $b_1 + \cdots +b_n =1$ and $b_i b_j =0$ for $i \neq j$. 
  Projections are given by the standard basis of $B^n$ and composition is defined by matrix multiplication as above. 
\end{definition}

%As expected, a model of the theory $H_B$ is a $B$-module $X$ whose $n$-ary operations $f(x_1, \ldots, x_n) =b_1x_1 + \cdots + b_n x_n$ with $b_i \in B$ and $x_i \in X$ satisfy $b_1 + \cdots +b_n =1$ and $b_i b_j =0$ for $i \neq j$. 
Boolean rings can be equivalently described as Boolean algebras defining $a \land b := ab$, $a\lor b= a + (1-a)b$, $\lnot a:=1-a$. 
We will freely switch between the two descriptions, because some results are better expressed in terms of Boolean algebras.


Note that, 
%Let $a,b \in B$;
% be elements of a Boolean ring. 
if $a b = a \land b = 0$, then $a +b = a \lor b$. 
As a consequence, the operations in $H_B(n)$ are $(b_1, \ldots, b_n) \in B^n$ such that $b_1 \lor \cdots \lor b_n =1$ and $b_i  \land b_j =0$ for $i \neq j$. 

Observe that $H_B(0)=\varnothing$ unless $B$ is degenerate, that is $0 = 1$, in which case $H_B(0)=\{*\}$. 


\begin{lemma}
  \label{lem:bool-hyper}
  Let $B$ be a Boolean ring. 
  The algebraic theory $H_B$ is hyperaffine. 
\end{lemma}
  \begin{proof}
    As $B$ is commutative, $H_B$ is commutative.
    Idempotence is ensured by the condition $b_1 + \cdots +b_n =1$: for every $(b_1, \ldots, b_n) \in H_B(n)$ and $(a_1, \ldots, a_m)\in H_B(m)$
    \begin{equation*} 
    \begin{bmatrix}
      b_1 \\
      \vdots \\
      b_n 
      \end{bmatrix} \circ  
      \begin{bmatrix}
      a_1 & \cdots & a_1 \\
      \vdots & \ddots & \vdots \\
      a_m & \cdots & a_m
      \end{bmatrix} =
      \begin{bmatrix}
      a_1(b_1 + \cdots + b_n)\\
      \vdots \\
      a_m(b_1 + \cdots + b_n) 
      \end{bmatrix}= 
      \begin{bmatrix}
      a_1 \\
      \vdots \\
      a_m \\
      \end{bmatrix}
\end{equation*} 
    The fact that $b_i ^2= b_i$ for all $i$ and $b_i b_j =0$ for $i \neq j$ precisely yields that every operation splits: 
    \small
    \begin{equation*}
    \begin{bmatrix}
      b_1 \\
      \vdots \\
      b_n \\
      \end{bmatrix} \circ  
      \begin{bmatrix}
      a^1_{11}b_1 + \cdots + a^1_{n1}b_n & \cdots & a^n_{n1}b_1 + \cdots + a^n_{n1}b_n \\
      \vdots & \ddots & \vdots \\
      a^1_{1m}b_1 + \cdots + a^1_{nm}b_n & \cdots & a^n_{1m}b_1 + \cdots + a^n_{nm}b_n
      \end{bmatrix} =
      \begin{bmatrix}
      a^1_{11}b_1 + \cdots + a^n_{n1} b_n\\
      \vdots \\
      a^1_{1m}b_1 + \cdots + a^n_{nm} b_n 
      \end{bmatrix}
    \end{equation*}
    \normalsize
    for all $(b_1, \ldots, b_n) \in H_B(n)$.
\end{proof}

Recall that we identify $f \in T(n)$ with the element $f(x_1, \ldots, x_n)$ in the free model on $n$ generators. 

\begin{lemma} 
  \cite[Proposition 4.4]{G24}
  \label{lem:hyper-bool}
    Let $T$ be a hyperaffine theory.  
    Then the set $T(2)$ is a Boolean algebra with the operations defined as follows, for every $a(x,y), b(x,y) \in T(2)$:
    \begin{itemize}
        \item $(a \land b)(x,y):=a(b(x,y),y)$; 
        \item $(a \lor b)(x,y):=a(x,b(x,y))$; 
        \item $(\lnot a)(x,y):=a(y,x)$; 
        \item $0(x,y):=y$ and $1(x,y):=x$. 
    \end{itemize}
\end{lemma}

\begin{comment}
\begin{proof}
  The only nontrivial identity to prove is $(a \land b) \lor (a \land c) = a \land (b \lor c)$. 
      It can be obtained, through some calculations, as a consequence of Lemma \ref{lem:c5} in the form: 
      %\begin{equation*}
        $a(b(a(x,y),a(z,w)),c(a(x,y),a(z,w))) = a(b(x,z),c(y,w)) \text{.}$
      %\end{equation*}
\end{proof}
\end{comment}

When the Boolean ring $B$ is degenerate, we have that $H_B=U$, the degenerate theory. 
%We will see that there is bijective correspondence between hyperaffine theories different from $U'$ and Boolean rings. 
%This correspondence is imperfect in the sense that the other degenerate theory $U'$ does not come from a ring. 

Johnstone \cite[p.~461]{J90} introduced the notion of coefficient to prove that, if $H_n$ is the hyperaffine theory generated by a single operation of arity $n$, then the category of models of $H_n$ is equivalent to the product $\mathbf{Set}^n$.  
The notion of coefficient will play a central role in what follows. 

\begin{definition}
  \label{def:coord}
Let $T$ be an algebraic theory, and let $f \in T(n)$. 
We call the following elements of $T(2)$ \emph{coefficients} of $f$: 
\begin{equation*}
  f[1](x,y):=f(x,y,\ldots,y), \ldots, f[n](x,y):=f(y,\ldots,y,x)\text{.}
\end{equation*}
\end{definition}

The next proposition shows in particular the fact that any operation in a hyperaffine theory 
is completely determined by 
its coefficients.   

\begin{proposition}
  \label{prop:coord2}
  Let $T$ be a hyperaffine theory and let $B$ be the Boolean ring $T(2)$.  
  The following hold: 
  \begin{enumerate}
    \item if $f \in T(n)$, then $f[1] \lor \cdots \lor f[n] =1$ and $f[i]f[j]=0$ for $i \neq j$; 
    \item for $f,g \in T(n)$, if $f[i]=g[i]$ for all $1 \le i \le n$, then $f=g$; 
    \item for every $(b_1, \ldots, b_n)$ such that $b_1 \lor \cdots \lor b_n =1$ and $b_i b_j =0$ for $i \neq j$, there is $f \in T(n)$ such that $f[i]=b_i$ for all $1 \le i \le n$. 
  \end{enumerate}
\end{proposition}
\begin{proof}
  Regarding the first item, the fact that $f[i] f[j]=0$ for each $i \neq j$ follows applying idempotence and \eqref{eq:c5} of Lemma \ref{lem:c5}; 
    that $f[1] \lor \cdots \lor f[n]=1$
    follows from associating the disjunction to the right and repeatedly applying \eqref{eq:c5}.  
  The second item is \cite[Lemma 4.6(i)]{G24}. 
  Finally, we prove the third item. 
  The proof is by induction on $n$ (cf. the proof of \cite[Proposition 4.7]{G24}).  
  If $n=1$, then $f(x)=x$. 
  Assume that the result holds for $n-1$, and let $(b_1, \ldots, b_n)$ be such that $b_1 \lor \cdots \lor b_n =1$ and $b_i \land b_j =0$ for $i \neq j$.   
  By inductive assumption there is $g \in T(n-1)$ such that $g[1]=b_1 \lor b_2$ and $g[i]=b_{i+1}$ for $2 \le i \le n$.
  Let $f(x_1, \ldots, x_n):=b_1(x_1, g(x_2, \ldots, x_n))$. 
  We prove that $f[i]=b_i$; we separate the case $i=1$, which does not depend on the inductive assumption,
  \begin{align*}
    f[1](x,y) & = b_1(x,g(y,\ldots,y)) & \\
    & = b_1(x,y) & g \text{ is idempotent} 
  \end{align*}
  from the case $i\neq1$:
  \begin{align*}
    f[i](x,y) & = b_1(y,g[i-1](x,y)) & \\
    & = \lnot b_1 (g[i-1](x,y),y) & \\
    & = (\lnot b_1 \land g[i-1])(x,y) &  \\
    & = (\lnot b_1 \land b_i)(x,y) & \text{by inductive assumption}
  \end{align*}
  for all $2 \le i \le n$, but $\lnot b_1 \land b_i =b_i$ as $b_1 \lor \cdots \lor b_n =1$ and $b_i \land  b_j =0$ for $i \neq j$. 
  This concludes the proof. 
%  The proof is similar to that of Proposition \ref{prop:coord}. 
%  It is therefore omitted here and deferred to the Appendix. 
%  See also \cite[Lemma 4.6]{G24}. 
\end{proof}

This leads us to a new proof of one of the results of \cite[Theorem 4.9]{G24}.

\begin{lemma}
  \label{lem:iso}
  Let $B$ be a Boolean ring. 
  The ring $H_B(2)$ of Lemma \ref{lem:hyper-bool} 
  %with Mal'cev operation $p=(1,-1,1) \in A_R(3)$ 
  is isomorphic to $B$. 
\end{lemma}
\begin{proof}
    Recall that $H_B(n)=\{(b_1, \ldots, b_n) \in B^n : b_1 \lor \cdots \lor b_n =1 \text{ and } b_i \land b_j =0 \text{ for } i \neq j\}$. 
    In particular $H_B(2) =\{(b,\lnot b) : b \in B\}$. 
    Then, by definition: 
    \begin{align*}
      (a, 1 -a) \lor (b,1- b) & = 
      \begin{bmatrix}
      1 & b \\
      0 & 1-b
      \end{bmatrix} 
      \begin{bmatrix}
      a \\
      1-a 
      \end{bmatrix} 
      %= (a+b+ab, 1-(a+b+ab)) 
      = (a \lor b, 1-(a \lor b)) \\
      (a,1-b) \land (b,1-b) & = 
      \begin{bmatrix}
      b & 0 \\
      1-b & 1 
      \end{bmatrix} 
      \begin{bmatrix}
      a \\
      1-a 
      \end{bmatrix} = 
      (a \land b, 1-(a \land b)) \\
      \lnot (b,1-b) & = 
      \begin{bmatrix}
      0 & 1\\
      1 & 0
      \end{bmatrix} 
      \begin{bmatrix}
      b \\
      1-b
      \end{bmatrix} = 
      (\lnot b, b) 
    \end{align*}
    This concludes the proof. 
\end{proof}

\begin{theorem}
  \label{thm:equibool}
  The full subcategory of hyperaffine theories different from $U'$ is equivalent to the category of Boolean rings. 
\end{theorem}
\begin{proof}
    %The proof is almost identical to that of Theorem \ref{thm:equibool}. 
    Let $H \neq U'$ be a hyperaffine theory, and let $B$ be a Boolean ring. 
    The two assignments $H \mapsto H(2)$ and $B \mapsto H_B$ are functorial. 
    By Lemma \ref{lem:iso} $B$ is isomorphic to $H_B(2)$. 
    For each $n \in \mathbb{N}$, we consider $\varphi: H(n) \to H(2)^n$ given by $\varphi(f) = (f[1], \ldots, f[n])$. 
    By Proposition \ref{prop:coord2} (1) $\varphi$ is well-defined, and by Proposition \ref{prop:coord2} (2) and (3) $\varphi$ is a bijection onto the image $\{(b_1, \ldots, b_n) \in H(2)^n : b_1 \lor \cdots \lor b_n =1, b_i b_j =0 \}$. 
    We prove that $\varphi$ is a morphism of theories, i.e. that for every $f \in H(n)$ and $g_1, \ldots, g_n \in H(m)$
  \begin{equation*}
    \varphi(f)(\varphi(g_1), \ldots, \varphi(g_n))(x_1, \ldots, x_m) = \varphi(f(g_1, \ldots,g_n))(x_1, \ldots, x_m) \text{.} 
  \end{equation*}
  By definition of $\varphi$ this amounts to prove that 
  \begin{equation*}
    \begin{bmatrix}
      g_1[1] & \cdots & g_n[1] \\
      \vdots & \ddots & \vdots \\
      g_1[m] & \cdots & g_n[m] 
    \end{bmatrix}
    \begin{bmatrix}
      f[1] \\
      \vdots \\
      f[n]
    \end{bmatrix}
    (x,y) = 
    \begin{bmatrix}
    f(g_1, \ldots, g_n)[1] \\
    \vdots \\
    f(g_1, \ldots, g_n)[m]
    \end{bmatrix}
    (x,y) \text{,}
  \end{equation*}
  i.e. that for every $i=1, \ldots, m$, 
  \begin{equation*}
    f[1]g_1[i] (x,y) \lor \cdots  \lor f[n]g_n[i](x,y)  = f(g_1[i](x,y), \ldots, g_n[i](x,y)) \text{.}
  \end{equation*} 
  For simplicity's sake, fix $i=1$. 
  Employing idempotence of $f$ and \eqref{eq:c5} Lemma \ref{lem:c5}, one can see that 
  \begin{align*}
    f[1]g_1[1] (x,y) \lor f[2]g_2[1](x,y) = f(g_1[1](x,y), g_2[1](x,y), y, \ldots, y)
  \end{align*}
  and then, inductively, 
  $$f[1]g_1[1] (x,y) \lor \cdots \lor f[n]g_n[1](x,y)=f(g_1[1](x,y), \ldots, g_n[1](x,y))\text{.}$$ 
  As the calculation with $i=2, \ldots, n$ is similar, we obtain the desired result. 
\end{proof}

\subsection{Commutative rings and affine theories}
%dire che i moduli sono algebra di questa teoria. 
%We omit precise definitions but roughly a semiring R is a ‘ring without subtraction’ and an R-module is an (additive) commutative monoid with a compatible (left) action of R. 
%Modules are the algebras of a theory TR, where TR(n) = Rn identifying (r1, . . . , rn) in Rn with r1x1 + ... + rnxn [5, 4]. 
%Switching back to monads briefly, TR (X ) is the set of finitely supported functions X → R.
In the case of affine theories, 
the assumption that every operation is a decomposition operation (cf. Definition \ref{def:icd}) is exchanged for 
the requirement of the existence of a Mal'cev operation. 
Note that the only model of a Mal'cev operation $p$ that splits is degenerate: 
\begin{equation*}
  y=p(y,x,x)=p(p(x,x,y), p(x,y,y), p(x,y,y)) = p(x,y,y) = x \text{,}
\end{equation*}
and therefore there are only two theories that are affine and hyperaffine at the same time: the degenerate ones. 

\begin{definition}
  \label{def:ar}
  For a given commutative ring $R$, we define the algebraic theory $A_R$ whose operations in $A_R(n)$ are $(r_1, \ldots, r_n) \in R^n$ such that $r_1 + \cdots +r_n =1$. 
\end{definition}

%Clearly, a model of the theory $A_R$ is an $R$-module $X$ whose $n$-ary operations $f(x_1, \ldots, x_n) =r_1x_1 + \cdots + r_n x_n$ with $r_i \in R$ and $x_i \in X$ satisfy $r_1 + \cdots +r_n =1$. 

\begin{lemma}
  The algebraic theory $A_R$ is idempotent and commutative. 
  Moreover, the operation $(1,-1,1) \in A_R(3)$ is a Mal'cev operation. 
  %In particular, $\mathbb{A}_R$ is affine according to Definition \ref{def:hyper-affine}. 
\end{lemma}
 \begin{proof}
    The proof is similar to that of Lemma \ref{lem:bool-hyper}. 
    The commutativity of the monoid $(R, \cdot, 1)$ ensures that $A_R$ is commutative, and the condition $r_1 + \cdots + r_n =1$ that $A_R$ is idempotent.
    Finally, $(1,-1,1)$ is a Mal'cev operation: 
    \begin{equation*} 
    \begin{bmatrix}
      1 \\
      -1\\
      1 
      \end{bmatrix} \circ  
      \begin{bmatrix}
      a_1 & b_1 & b_1 \\
      \vdots & \vdots & \vdots \\
      a_m & b_m & b_m
      \end{bmatrix} =
      \begin{bmatrix}
      a_1\\
      \vdots \\
      a_m 
      \end{bmatrix}
\end{equation*} 
for any $(a_1, \ldots, a_m), (b_1, \ldots, b_m) \in T(m)$ and similarly for the other condition. 
\end{proof}

%The following lemma was proved in \cite[Theorem 6.3.3]{RS02}. 
We prove the analogous of Lemma \ref{lem:hyper-bool}. 
See \cite[Theorem 6.3.3]{RS02} for the similar result stating that a free Mal'cev mode on two generators can be endowed with the structure of commutative ring. 

\begin{lemma}
  \label{lem:ring}
    If $T$ is affine, then $T(2)$ is a commutative ring with the operations defined as follows, for every $a(x,y), b(x,y) \in T(2)$: 
  \begin{itemize}
    \item $(a + b)(x,y) := p(a(x,y),y,b(x,y))$; 
    \item $(a \cdot b)(x,y) := a(b(x,y),y)$; 
    \item $-a(x,y):=p(y,a(x,y),y)$; 
    \item $0(x,y):=y$ and $1(x,y):=x$ 
  \end{itemize} 
\end{lemma}
\begin{proof}
    Associativity and commutativity of the addition follow from Lemma \ref{lem:asscomm} (M1) and (M2). 
    Moreover, 
    \begin{align*}
      (a + 0)(x,y) & = p(a(x,y), y, y) & \\
      & = a(x,y) & p \text{ is Mal'cev} 
    \end{align*}
    \begin{align*}
      (a + (-a))(x,y) & = p(a(x,y),y,p(y,a(x,y),y)) & \\
      & = p(p(a(x,y),y,y),a(x,y),y) & \text{by Lemma \ref{lem:asscomm} (M1)}  \\
      & = p(a(x,y), a(x,y), y) & p \text{ is Mal'cev} \\ 
      & = y & p \text{ is Mal'cev} \\
      & = 0(x,y) \text{.} 
    \end{align*}
    Associativity and unitality of the multiplication are immediate. 
    We prove commutativity: 
    \begin{align*}
      (a \cdot b)(x,y) & = a(b(x,y), y) & \\
      & = a(b(x,y), b(y,y)) & b \text{ is idempotent}  \\ 
      & = b(a(x,y), a(y,y)) & a,b \text{ commute} \\
      & = b(a(x,y), y) & a \text{ is idempotent} \\
      & = (b\cdot a)(x,y)  \text{.} 
    \end{align*}
    Finally, 
    \begin{align*}
      r(a+b)(x,y) & = r(p(a(x,y),y,b(x,y)),y) & \\
      & = r(p(a(x,y),y,b(x,y)),p(y,y,y)) & p \text{ is idempotent}  \\ 
      & = p(r(a(x,y),y), r(y,y), r(b(x,y),y)) & p,r \text{ commute}  \\ 
      & = p(r(a(x,y),y), y, r(b(x,y),y)) & r \text{ is idempotent} \\
      & = (ra +rb)(x,y) \text{.} 
    \end{align*}
    This completes the proof. 
  \end{proof}

\begin{lemma}
  \label{lem:pl}
  If $T$ is an affine theory, then 
  \begin{equation*}
    p(a(x,y), b(x,y), c(x,y))=(a-b+c)(x,y)\text{.}
  \end{equation*}
\end{lemma}
\begin{proof}
  Let $p$ be the Mal'cev operation. Then, exploiting Lemma \ref{lem:asscomm}
  \begin{align*}
    p(a(x,y),b(x,y),c(x,y)) & = p(a(x,y),b(x,y),p(y,y,c(x,y))) & \\
    & = p(p(a(x,y),b(x,y),y),y,c(x,y)) & \text{(M1)} \\
    & = p(p(p(a(x,y),y,y),b(x,y),y),y,c(x,y)) & \\
    & = p(p(a(x,y),y,p(y,b(x,y),y)),y,c(x,y)) & \text{(M1)}\\
    & = p((a-b)(x,y),y,c(x,y)) & \\
    & = (a-b+c)(x,y) & 
  \end{align*}
  for every $a,b,c \in T(2)$. 
\end{proof}

%Starting from a ring, applying both constructions returns an isomorphic ring.

\begin{lemma}
  \label{lem:unit}
  Let $R$ be a commutative ring. 
  The ring $A_R(2)$ of Lemma \ref{lem:ring} 
  %with Mal'cev operation $p=(1,-1,1) \in A_R(3)$ 
  is isomorphic to $R$. 
\end{lemma}
\begin{proof}
    Recall that $A_R(n)=\{(r_1, \ldots, r_n) \in R^n : r_1 + \cdots +r_n =1\}$. 
    In particular $A_R(2) =\{(r,1-r) : r \in R\}$. 
    Then, by definition: 
    \begin{align*}
      (r,1-r) + (s,1-s) & = 
      \begin{bmatrix}
      r & 0 & s\\
      1-r & 1 & 1-s
      \end{bmatrix} 
      \begin{bmatrix}
      1 \\
      -1 \\
      1
      \end{bmatrix} = 
      (r+s, 1-r-s) \\
      (r,1-r) \cdot (s,1-s) & = 
      \begin{bmatrix}
      s & 0 \\
      1-s & 1 
      \end{bmatrix} 
      \begin{bmatrix}
      r \\
      1-r 
      \end{bmatrix} = 
      (rs, 1-rs) \\
      -(r,1-r) & = 
      \begin{bmatrix}
      0 & r & 0\\
      1 & 1-r & 1
      \end{bmatrix} 
      \begin{bmatrix}
      1 \\
      -1 \\
      1
      \end{bmatrix} = 
      (-r, 1+r) 
    \end{align*}
    This concludes the proof. 
\end{proof}

Now, we prove the analogue of Proposition \ref{prop:coord2}.

\begin{proposition}
  \label{prop:coord}
  Let $T$ be an affine theory.
  Let $R$ be the commutative ring $T(2)$. 
  The following hold: 
  \begin{enumerate}
    \item if $f \in T(n)$, then $f[1]+\cdots+f[n]=1$ in $R$; 
    \item for $f,g \in T(n)$, if $f[i]=g[i]$ for all $1 \le i \le n$, then $f=g$; 
    \item for every $(r_1, \ldots, r_n) \in R^n$ such that $r_1+\cdots+r_n=1$, there is $f \in T(n)$ such that $f[i]=r_i$ for all $1 \le i \le n$. 
  \end{enumerate}
\end{proposition}
\begin{proof}
    Regarding the first item, we prove that $f[1] + \cdots + f[n] = 1$. 
    Observe that 
    \begin{align*}
      (f[1]+f[2])(x,y) & = p(f[1](x,y),y,f[2](x,y)) & \\
      & = p(f(x,y, \ldots, y) , f(y,y \ldots, y), f(y,x, \ldots, y)) & f \text{ is idemp.}\\
      & = f(p(x,y,y), p(y,y,x), \ldots, p(y,y,y)) & p,f \text{ commute} \\
      & = f(x,x,y \ldots, y) & p \text{ is Mal'cev} 
    \end{align*}
    and, therefore, inductively, 
    \begin{equation*}
      (f[1]+ \cdots + f[n])(x,y) = f(x,x, \ldots,x) =x \text{.}
    \end{equation*}

    For the second item, observe that 
    \begin{align*}
      p(f[1](x_1,y),y,f[2](x_2,y)) & = p(f(x_1,y, \ldots, y) , f(y,y \ldots, y), f(y,x_2, \ldots, y)) & \\
      & = f(p(x_1,y,y), p(y,y,x_2), \ldots, p(y,y,y)) &  \\
      & = f(x_1,x_2,y \ldots, y) &
    \end{align*}
    and, therefore, inductively, 
    \begin{equation*}
      p(f(x_1,\ldots,x_{n-1},y), y,f[n](x_n,y)) = f(x_1,x_2, \ldots,x_n)\text{.}
    \end{equation*}
    This implies that if $f[i]=g[i]$ for all $i=1, \ldots, n$, then $f = g$. 

  Finally, we prove the third item. 
  %First of all, observe that, if $(r_1,r_2,r_3)=(1,-1,1)$, then any Mal'cev operation $p \in T(3)$ is such that $p[i]=r_i$ for $i=1,2,3$.  
  The proof is by induction on $n$. 
  If $n=1$, then clearly $f(x)=x$. 
  Assume that the result holds for $n-1$, and let $(r_1, \ldots, r_n) \in R^n$ be such that $r_1 + \cdots +r_n=1$. 
  By inductive assumption there is $g \in T(n-1)$ such that $g[1]=r_1 + r_2$ and $g[i]=r_{i+1}$ for $2 \le i \le n$.
  Let $f(x_1, \ldots, x_n):=p(r_2(x_2,x_1),x_1,g(x_1,x_3,\ldots,x_n))$. 
  We prove that $f[i]=r_i$; we separate three cases $i=1,2$ and $3 \le i \le n$:  
  \begin{align*}
    f(x,y,\ldots,y) & = p(r_2(y,x),x,g(x,y,\ldots,y)) & \\
    & = r_2(y,x) -x +g[1](x,y) & \text{ by Lemma \ref{lem:pl}}\\
    & = (1-r_2 -1+r_1+r_2)(x,y) & \text{ by inductive assumption} \\
    & = r_1(x,y)
  \end{align*} 
  \begin{align*} 
    f(y,x,y,\ldots,y) & = p(r_2(x,y),y,g(y, \ldots, y)) & \\
    & = p(r_2(x,y),y,y) & g \text{ is idempotent}\\
    & = r_2(x,y) & p \text{ is Mal'cev} 
  \end{align*}
  \begin{align*}
    f[i](x,y) & = p(r_2(y,y),y,g[i-1](x,y)) & 3 \le i \le n \\
    & = p(y,y,g[i-1](x,y)) & r_2 \text{ is idempotent} \\
    & = g[i-1](x,y) & p \text{ is Mal'cev} \\
    & = r_i(x,y) & \text{ by inductive assumption} 
  \end{align*}
  This concludes the proof. 
\end{proof}

We have all the ingredients to present one of the main theorems. 

\begin{theorem}
  \label{thm:equiring}
  The full subcategory of affine theories different from $U'$ is equivalent to the category of commutative rings. 
\end{theorem}
\begin{proof}
  Let $A\neq U'$ be an affine theory, and let $R$ be  a commutative ring. 
  The two assignments $A \mapsto A(2)$ and $R \mapsto A_R$ are functorial. 
  Lemma \ref{lem:unit} says that $R$ is isomorphic to $A_R(2)$. 
  For each $n \in \mathbb{N}$, consider the function $\varphi: A(n) \to A(2)^n$ such that $\varphi(f) = (f[1], \ldots, f[n])$. 
  By Proposition \ref{prop:coord}, $\varphi$ is a bijection onto the image $\{(r_1, \ldots, r_n) \in A(2)^n : r_1 + \cdots + r_n =1\}$. 
  We prove that $\varphi$ is a morphism of theories, i.e. that for every $f \in A(n)$ and $g_1, \ldots, g_n \in A(m)$
  \begin{equation*}
    \varphi(f)(\varphi(g_1), \ldots, \varphi(g_n))(x_1, \ldots, x_m) = \varphi(f(g_1, \ldots,g_n))(x_1, \ldots, x_m) \text{.} 
  \end{equation*}
  By definition of $\varphi$ this amounts to prove that 
  \begin{equation*}
    \begin{bmatrix}
      g_1[1] & \cdots & g_n[1] \\
      \vdots & \ddots & \vdots \\
      g_1[m] & \cdots & g_n[m] 
    \end{bmatrix}
    \begin{bmatrix}
      f[1] \\
      \vdots \\
      f[n]
    \end{bmatrix}
    (x,y) = 
    \begin{bmatrix}
    f(g_1, \ldots, g_n)[1] \\
    \vdots \\
    f(g_1, \ldots, g_n)[m]
    \end{bmatrix}
    (x,y) \text{,}
  \end{equation*}
  i.e. that for every $i=1, \ldots, m$, 
  \begin{equation*}
    f[1]g_1[i] (x,y) + \cdots +f[n]g_n[i](x,y)=f(g_1[i](x,y), \ldots, g_n[i](x,y)) \text{.}
  \end{equation*} 
  For simplicity's sake, fix $i=1$. 
  Observe that 
  \begin{align*}
    f[1]g_1[1] (x,y) + f[2]g_2[1](x,y) & = p(f[1](g_1[1](x,y),y),y, f[2](g_2[1](x,y),y)) \\
    & = p(f[1](g_1[1](x,y),y),f(y,\ldots,y), f[2](g_2[1](x,y),y))  \\
    & = f(p(g_1[1](x,y),y,y),p(y,y,g_2[1](x,y)), y, \ldots, y) \\
    & = f(g_1[1](x,y), g_2[1](x,y), y, \ldots, y)
  \end{align*}
  and inductively, $f[1]g_1[1] (x,y) + \cdots +f[n]g_n[1](x,y)=f(g_1[1](x,y), \ldots, g_n[1](x,y))$ as desired.
  As the calculation with $i=2, \ldots, n$ is similar, we obtain the desired result. 
\end{proof}

\section{Hyperaffine theories and algebraic logic.}
\label{sec:nba}
  In \cite{Dicker} it is shown that the variety of Boolean algebras can be axiomatised using two nullary operations, $0$ and $1$, and a ternary operation $q$ called `conditional disjunction'. 
Expressed in terms of the Boolean operators, conditional disjunction is $q(a,x,y)=(a \land x) \lor (\lnot a \land y)$.
% i.e. if $a$ then $x$ else $y$.  
  The restriction of the conditions of Definition \ref{def:hyper-affine} defining hyperaffine theories to binary operations yields the following axiomatisation of Boolean algebras:
  % in the signature with a ternary operation $q$ and two constants $0,1$: 
\begin{itemize}
  \item $q(a,1,0)=a$; \hfill %cf. \eqref{eq:c12} 
  \item $q(1,a,b)=a$ and $q(0,a,b)=b$; %\hfill cf. \eqref{eq:c12} 
  \item $q(a, q(b,x,y), q(c,x,y))= q(q(a,b,c),x,y)$; %\hfill cf. \eqref{eq:c3}
  \item $q(a,x,x)=x$; %\hfill cf. Definition \ref{def:hyper-affine}
  \item $q(a, q(b,x,y), q(b,z,w)) = q(b, q(a,x,z), q(a,y,w))$; %\hfill cf. Definition \ref{def:hyper-affine}
  \item $q(a, q(a,x,y), q(a,z,w)) = q(a,x,w)$. %\hfill cf. Definition \ref{def:hyper-affine} 
\end{itemize}
In \cite{BS,SBLP20,BLPS18} these axioms have been adapted from dimension $2$ to
%(featuring two constants, $0$ and $1$, and a ternary operation $q$) 
dimension $n$, with $n$ constants $e_1, \ldots, e_{n}$ acting as projections and a generalised $(n+1)$-ary if-then-else $q$. 

\begin{definition} \cite[Definition 3.7]{BS}
  A \emph{Boolean algebra} of \emph{dimension $n$} ($n$BA) is a set $A$ endowed with nullary operations $e_1, \ldots, e_n$ and a $(n+1)$-ary operation $q : A^{n+1} \to A$ satisfying the following identities:
  \begin{description}
  \item[H1] $q(e_i,x_1,\ldots,x_n) = x_i$ for every $i = 1,\ldots,n$;
  \item[H2] $q(y,x,\ldots,x) = x$; 
  \item[H3] $q(y,q(y,x^1_1,x^1_2,...,x^1_n),\ldots,q(y,x^n_1,x^n_2,\ldots,x^n_n)) = q(y,x^1_1,\ldots,x^n_n)$; 
  \item[H4] $q(y,q(z_1,x^1_1, \ldots,x^1_n),\ldots,q(z_n,x^n_1, \ldots,x^n_n)) \\
  = q(q(y,z_1, \ldots, z_n),q(y, x^1_1,\ldots,x^n_1),\ldots,q(y,x^1_n,\ldots,x^n_n))$;
  \item[H5] $q(y,e_1,...,e_n) = y$.
\end{description}
\end{definition}

If $A$ is an $n$BA, let $t(a,a_1, a_2) := q(a,a_1, a_2, e_3, \ldots, e_n)$ for every $a,a_1,a_2 \in A$. 
The set $B_A:=\{a \in A : t(a,b,b)=b\}$ is a Boolean algebra with respect to the ternary operation $t$ and constants $e_1, e_2$; 
see \cite[Lemma 4]{BLPS18} and subsequent discussion. 
The Boolean algebra $B_A$ is called the \emph{Boolean algebra of the coordinates} of $A$. 
The \emph{coordinates} of an element $a \in A$ are the elements of $B_A$ defined as 
\begin{equation*}
  a[1]=q(a,e_1,e_2, \ldots, e_2), \ldots, a[n]=q(a,e_2, \ldots,e_2,e_1)\text{.}
\end{equation*} 
See \cite[Definition 10, Lemma 6]{BLPS18}.   
This is similar to the Definition \ref{def:coord} of coefficients given above. 
%The xxx is a Boolean algebra and is called the Boolan algebra of the coordinates 
%The axiomatisation presented in \cite[Definition~14]{SBLP20} is equivalent to the $(n+1)$-ary form of axioms (H1)-(H6), thanks to Lemma~\ref{lem:c5}.

The connection with hyperaffine theories is the following.
\begin{theorem}
  \label{thm:nba}
  If $T$ is a hyperaffine theory, then $T(n)$ is an $n$BA by letting 
  $e_i:=\pi^n_i$ for $i =1,\ldots, n$ and by defining 
  \begin{equation*}
    q: T(n)^{n+1} \to T(n) \quad \text{ as } \quad \circ: T(n) \times T(n)^n \to T(n)\text{.}
  \end{equation*}
Conversely, for any sequence $(A_n, q_n, e^n_1, \ldots, e^n_n)_{n \ge 2}$ of $n$BAs such that $B_{A_n}=B_{A_m}$ for every $n,m$
there is, up to isomorphism of theories, a unique hyperaffine theory $T$ such that $T(n)=A_n$,  
  \begin{equation*}
    q_n(a, b^1, \ldots, b^n) = a \circ (b^1, \ldots, b^n) \quad \text{ and } \quad e^n_i = \pi^n_i
  \end{equation*}
  for every $a, b^1, \ldots, b^n \in A_n$ and every $i=1, \ldots, n$. 
%the same Boolean algebra of coordinates uniquely determines a hyperaffine theory. 
\end{theorem}
%\end{comment}
\begin{proof}
  Let $T$ be hyperaffine and let $n \ge 1$.
  Then $T(n)$ satisfies (H1) by the definition of algebraic theory, more precisely condition \eqref{eq:c12}.  
  Identity (H2) follows from the fact that $T$ is idempotent, and (H3) because every operation splits. 
  Finally, $T(n)$ satisfies (H4) by Lemma \ref{lem:c5} and (H5) again by \eqref{eq:c12}. 

  Conversely, let $(A_n)_{n \ge 2}$ be a sequence of $n$BAs having the same Boolean algebra $B$ of the coordinates. 
  To avoid trivialities, we assume that $B$ is nondegenerate. 
  Then we let $T$ to be the hyperaffine theory $H_B$. 
  Consider the function $\psi: A_n \to H_B(n)$ mapping an element $a$ to its coordinates $\psi(a)=(a[1], \ldots, a[n])$. 
  By \cite[Lemma 8 (i)]{BLPS18} $\psi$ is well-defined. Moreover, it is injective by \cite[Lemma 9]{BLPS18} and a surjective homomorphism by \cite[Theorem 8]{BLPS18}. 
  The fact that $\pi^n_i = e^n_i$ for every $i=1, \ldots, n$ is immediate and that $q_n(a, b^1, \ldots, b^n) = a \circ (b^1, \ldots, b^n)$
  for every $a, b^1, \ldots, b^n \in A_n$ follows from \cite[Lemma 8(ii)]{BLPS18}.
\end{proof}


On the one hand, $n$BAs have inspired work aimed at generalising the proof theory of classical propositional logic to the case of an arbitrary finite number of symmetric truth values \cite{BCS24}.
On the other, they have been discovered to have a connection with the skew Boolean algebras introduced by Leech \cite{Leech1,Leech2,Leech3}. 
More precisely, it was shown in \cite{BS} that 
any $n$BA $A$ contains a skew cluster of isomorphic right-handed skew Boolean algebras, called its \emph{skew reducts};
moreover, the skew reducts of $A$ allow to recover $A$ perfectly. 
By exploiting the correspondence between affine theories and commutative rings from Theorem \ref{thm:equiring}, we aim to define an $n$-dimensional commutative ring via 
%the Mal'cev operation and 
axioms similar to (H1)-(H5). 
It is reasonable to expect that this approach may yield a definition of a skew ring and a result similar to the one established in \cite{BS} and mentioned previously. 
We defer this investigation to future research.

\section{Boolean actions} 
\label{sec:spaces}
%In this section, we take a closer look at the models of hyperaffine theories, and affine theories whose ring of coefficients is Boolean. 
Given a Boolean ring $B$, there are two ways to encode $B$ as an algebraic theory: 
either as the hyperaffine theory $H_B$ or as the affine theory $A_B$. 
Models of $H_B$ are the well-known $B$-sets \cite{B91}. 
It is therefore of interest to investigate models of an affine theory when the coefficient ring is Boolean, a task we undertake in this section. 
We observe that the models of $H_B$ highlight the set-theoretic properties of $B$, 
whereas the models of $A_B$ emphasise its algebraic structure 
(cf. Theorem 5.7 and the role of Boolean vector spaces).
%with the goal of comparing the two. 
%We fix a Boolean ring $B$. 
%We recall that models of $A_B$ (Definition \ref{def:ar}) are typically called \emph{affine} $B$-modules. 
%We shall prove that models of the former are generalisations of affine spaces over a field, while models of the latter conform to a well-known definition of Boolean action due to Bergman \cite{B91}. 
We begin with the following observation. 
%Observe that Definition \ref{def:affspace} does not make evident the fact that they form an algebraic category. 
%The starting point is the observation that an affine theory is determined by the set of its binary operations, which form a ring (Lemma \ref{lem:ring}).
%Moreover, idempotence and commutativity are essential requirements.
%Given a commutative ring $R$, the following conditions on a model $X$ of $\mathbb{A}_R$ are therefore natural. 

\begin{lemma}
  \label{prop:nec}
  Let $B$ be a Boolean ring. 
  If $X$ is a model of $A_B$ or a model of $H_B$, then $X$ is endowed with a binary action of $B$, i.e. a function $B\times X^2 \to X$, satisfying the following axioms:
\begin{description}
  \item[R1] $b(x,x)=x$; 
  \item[R2] $b(c(x,y), c(z,w)) = c(b(x,z),b(y,w))$; 
  \item[R3] $1(x,y) =x$ and $0(x,y)=y$; 
  \item[R4] $a(b(x,y),c(x,y)) = (ab+(1-a)c)(x,y)$   
\end{description}
for all $x,y,z,w \in X$ and $a,b,c \in R$. 
\end{lemma}
\begin{proof}
  By Lemma \ref{lem:unit} and Lemma \ref{lem:iso}, $A_B(2)$ and $H_B(2)$ are Boolean rings isomorphic to $B$. 
  Note that (R1) follows from the idempotence of $A_B$ (or $H_B$), (R2) from the commutativity, while (R3) and (R4) from the fact that $X$ is a model (cf. conditions \eqref{eq:alpha1} and \eqref{eq:alpha2} of Section 2). 
\end{proof}
 
%We will see in Theorem \ref{thm:ringaction} when the conditions (R1)-(R4), necessary by Proposition \ref{prop:nec}, are also sufficient. 
%We will make formal these considerations shortly. 

%We observe that a model of $A_B$ (or $H_B$) can be empty. 
A set of axioms similar to (R1)-(R4) was given by Bergman with the purpose of capturing the concept of action of $B$ on a set.
%(the interested reader can find more about that in Section \ref{sec:more} of the Appendix). 
% When the ring $R$ is Boolean, a smiliar set of axioms was given by Bergman. 

\begin{definition} \cite[Definition 7]{B91}
  \label{def:bset}
  Let $B$ be a Boolean ring. 
  A \emph{$B$-set} is a set $X$ together with a function $b: X^2 \to X$ for each $b \in B$ satisfying the following requirements: 
\begin{description}
  \item[B1] $b(x,x)=x$; 
  \item[B2] $b(b(x,y), z) = b(x,z)= b(x,b(y,z))$; 
  \item[B3] $(1-b)(x,y)=b(y,x)$;    
  \item[B4] $0(x,y) =y$;
  \item[B5] $b(c(x,y),y)=(bc)(x,y)$  
\end{description}
for all $x,y \in X$ and $b,c \in B$. 
\end{definition}

\begin{remark}
  \label{rem:balgebra}
  More generally, an action of $B$ on an algebra is an action on the underlying set of the algebra, such that each operation $b \in B$ commutes with each operation of the algebra.
\end{remark}

We establish that (B1)-(B5) can be rewritten as (R1)-(R4) plus an additional axiom. 

\begin{lemma}
  \label{lem:rtob}
    For a Boolean ring $B$, \emph{(B1)-(B5)} are equivalent to \emph{(R1)-(R4)} plus
    \begin{description}
        \item[R5] $b(b(x,y), b(z,w)) = b(x,w)$
    \end{description} 
    for all $x,y,z,w \in X$ and $b\in B$. 
\end{lemma}
    \begin{proof}
        Assume (R1)-(R5). 
        (R5) combined with (R1) immediately gives (B2).  
        (R4) and (R3) give (B3). 
        (B5) follows from (R4) with $c=0$.   
        Conversely, assume (B1)-(B5). 
        The content of \cite[Proposition 11]{B91} is precisely that (B1)-(B5) imply (R2). 
        %From (B2) and (R5), (R2) also follows. 
        (R3) is obvious in light of (B3) and (B4). 
        It is not difficult to see that (R5) follows repeatedly applying (B2). 
        To obtain (R4), first observe that $ab+(1-a)c = ab \lor (1-a)c$ as $ab(1-a)c=0$. 
        Moreover, it easy to see that 
        \begin{equation}
          \label{eq:or}
          (a \lor b)(x,y) = a(x,b(x,y)) \text{.}
        \end{equation}
        Therefore  
        \begin{align*}
          (ab \lor (1-a)c)(x,y) & = a(b(x,a(y,c(x,y))),a(y,c(x,y))) & \eqref{eq:or}, \text{(B3), (B5)} \\
          & = a(a(b(x,y),b(x,c(x,y))),a(y,c(x,y))) & \text{ (R2)} \\
          & = a(b(x,y), c(x,y)) & \text{ (R5)}
        \end{align*}
        This concludes the proof. 
\end{proof}

Algebras for a hyperaffine theory $T$ are exactly sets equipped with the Boolean action of $T(2)$.
%We present a conceptually simpler (as it does not rely on Bergman's Theorem~\ref{thm:berg}) proof of 


\begin{corollary}
  \label{thm:models}
  Let $T$ be a hyperaffine theory, 
  and let $B:=T(2)$ be the Boolean ring associated with $T$. 
  Then the models of $T$ are sets equipped with an action of $B$ satisfying \emph{(R1)-(R5)}. 
  % or, equivalently, by Lemma \ref{lem:rtob}, $B$-sets according to Definition \ref{def:bset}. 
\end{corollary}
\begin{proof}
  By \cite[Proposition 3.16]{G24}, and by Lemma \ref{lem:rtob}. 
\end{proof}

\begin{comment}
\begin{proof} 
    The category of $B$-sets is clearly a variety, with algebraic theory $T$ generated by the binary operations $\{b(x,y): b \in B\}$ modulo the axioms (R1)-(R5). 
    The theory $T$ is idempotent by (R1), commutative by (R2), and every operation splits by (R5). 
    By Theorem \ref{thm:equibool} it is therefore of the form $H_{B'}$ for some Boolean ring $B'$, and we are done if $B \simeq B'$. 
    We know that $B' \simeq T(2)$. 
    Let $f(x,y)$ be an element of $T(2)$. 
    By (R3) and (B3) we can assume that $f(x,y)$ is the composition of binary operations of the form $d(x,y)$, with $d \in B$.
    We then apply (R4) in the form of the rewriting rule 
    \begin{equation*}
      a(b(x,y), c(x,y)) \to (ab + (1-a)c) (x,y)
    \end{equation*}
    to get a unique normal form $b(x,y)$ with $b \in B$. 
    This shows that $T(2) \simeq B$. 
\end{proof}
\end{comment}


Unlike the case of hyperaffine theories, binary operations do not \emph{completely} determine an affine theory, due to the presence of the ternary Mal'cev operation $p$. 
However, axioms (R1)-(R4) can describe models of an affine theory if the Mal'cev operation can be expressed as composition of binary operations, 
and this holds 
%For a ring $R$, the ternary operation $f(x,y,z)=x-y+z$ can be expressed as composition of binary ones 
iff there are no ring homomorphisms from $R$ to $\mathbb{F}_2$, where $\mathbb{F}_2$ is the field with two elements (see e.g. \cite{S77,IKS78}).  
In this case, models of an affine theory $T$ are sets with a binary action of the ring $R:=T(2)$ satisfying (R1)-(R4). 

Axioms in the general case can be deduced as follows. 

\begin{proposition} 
  \label{prop:6.3.4}
  For an affine theory $T$, and associated ring $R:=T(2)$, models of $T$ are sets equipped with binary operations $a: X^2 \to X$ for each $a \in R$ satisfying (R1)-(R4) and a ternary operation $p: X^3 \to X$ satisfying 
\begin{description}
  %\item[A0] $p$ is a Mal'cev operation; 
  \item[A1] $p$ is a Mal'cev operation;  
  \item[A2] $p$ commutes with itself and with every $a \in R$; 
  \item[A3] $p(y,x,y)=(1+1)(y,x)$; 
  \item[A4] $p(a(x,y), b(x,y), c(x,y)) = (a-b+c)(x,y)$.    
\end{description}
\end{proposition}
\begin{proof}
  A model of $T$ is an affine $R$-module, that is a $R$-module whose linear combinations $r_1x_1 + \cdots + r_n x_n$ satisfy $r_1 + \cdots + r_n =1$.
  Then the result follows from \cite[Theorem 6.3.4]{RS02}. 
\end{proof}

When $B$ is a Boolean ring, we obtain the following characterisation. 
\begin{theorem}
  \label{thm:new1}
  Let $B$ be a Boolean ring, and consider the affine theory $A_B$. 
  Models of $A_B$, together with the choice of an element $o \in X$, are precisely vector spaces over $\mathbb{F}_2$ equipped with an action of $B$ satisfying \emph{(R1)-(R4)} and such that 
  \begin{description}
  \item[L1] $b(x,y) + b(z,w) = b(x+z,y+w)$
  %\item[L2.] $a(x,y) + b(x,y) + c(x,y) = (a+b+c)(x,y)$
  \end{description} 
  for every $a,b \in B$ and $x,y,z,w \in X$.
\end{theorem}
\begin{proof}
  If $X$ models $A_B$, then $X$ is endowed with an action of $B$ and has a ternary operation satisfying (R1)-(R4) and (A1)-(A4).  
  We endow $X$ with the structure of $\mathbb{F}_2$-vector space as follows.
  %We assume that $X$ is nonempty.
  %The $B$-module structure on $X$ gives an Abelian group $(X,+,0)$, to which we add the action of $\mathbb{F}_2$: $1x=x$ and $0x=0$.
  We define $x+y:=p(x,o,y)$; 
  this sum is associative and commutative by (A1), (A2) and Lemma \ref{lem:asscomm}. 
  The element $o \in X$ is a zero for the sum by (A1). 
  Moreover, the inverse of $x$ is given by $x$ itself as (A3) shows: 
  \begin{equation*}
    x+x = p(x,o,x) = 0(x,o) =o\text{.}
  \end{equation*}
  Then, we add the action of $\mathbb{F}_2$: $1x:=x$, $0x:=o$. 
  The only nontrivial axiom is proven using (A3): 
  \begin{equation*}
    (1+1)x=0x=o=0(x,o)=(1+1)(x,o)=p(x,o,x)=x+x=1x+1x \text{.}
  \end{equation*}
  %We remark that $p(x,y,z)=x-y+z$, as two Mal'cev operations that commute with each other coincide (Lemma \ref{lem:asscomm}).
  Finally, using (A2) it is easy to derive (L1). 
  Conversely, given an $\mathbb{F}_2$-vector space $(X,+,0)$ with an action of $B$ that satisfies (R1)-(R4) and (L1), we define $p(x,y,z)=x+y+z$. 
  As $X$ has characteristic $2$, (A1) is verified. 
  (A3) follows by definition of $p$: 
  \begin{equation*}
    p(y,x,y)=y+x+y=(1+1)y+x=x=0(y,x)=(1+1)(y,x) \text{.}
  \end{equation*}
  Moreover, (A2) follows from (L1). 
  Now, observe that 
  \begin{align*}
    (a+b)(x,y) & = a(b(y,x),b(x,y)) & \text{(R4),(B3)} \\
    & = a(b(x,y)+b(y,x)+b(x,y),0+b(x,y)) & \\
    & = a(b(x,y)+b(y,x),0) + b(x,y) & \text{(L1),(R1)} \\
    & = a(b(x+y,y+x),0) + b(x,y) & \text{(L1)} \\
    & = a(b(x+y,x+y),0) + b(x,y) & \\
    & = a(x+y,0)+b(x,y) & \text{(L1),(R1)} \\
    & = a(x+y,y+y)+b(x,y) & \\
    & = a(x,y) + y + b(x,y) & \text{(L1),(R1)}
  \end{align*}
  and therefore 
  \begin{align*}
    (a+b+c)(x,y) & = a(b(c(x,y),c(y,x)),b(c(y,x),c(x,y))) & \text{(R4)} \\
    & = a(b(y,x)+x +c(y,x),b(x,y)+y+c(x,y)) & \\
    & = a(b(y,x),b(x,y)) +a(x,y) + a(c(y,x),c(x,y)) & \text{(L1)}\\
    & = a(x,y)+y+b(x,y) +a(x,y)+a(x,y)+y+c(x,y) & \\
    & = a(x,y) + b(x,y)+c(x,y) & 
    %& = p(a(x,y),b(x,y),c(x,y)) &\text{Lemma \ref{lem:asscomm}}
  \end{align*}
  %We leave to the reader to verify, using Lemma \ref{lem:asscomm}, that $a(x,y) - b(x,y)+c(x,y)=p(a(x,y),b(x,y),c(x,y)) 
  proving (A4). 
  % and (A4) from (L2), concluding the proof. 
\end{proof}

We sum up the results of this section in Table \ref{tab:axioms}.
In the table $B$ is a Boolean ring and $R$ is a commutative ring. 

\begin{table}%[h!]
  \caption{A summary of the results of this section}
    \label{tab:axioms}
\centering
\begin{tabular}{c c c}
\toprule
\textbf{Axioms} & \textbf{characterise} & \textbf{by virtue of}\\
\midrule
(B1)-(B5)  & models of $H_B$ & \cite[Proposition 3.16]{G24} \\
\midrule
(R1)-(R5)  & models of $H_B$ & Lemma \ref{lem:rtob} \\
\midrule 
(R1)-(R4), (A1)-(A4)  & models of $A_R$ & Proposition \ref{prop:6.3.4} \\
\midrule 
(R1)-(R4)  & models of $A_R$ with no $R \to \mathbb{F}_2$ & \cite{S77,IKS78} \\
\midrule 
(R1)-(R4), (L1) & models of $A_B$ & Theorem \ref{thm:new1} \\
\bottomrule
\end{tabular}
%\caption{Table to test captions and labels.}
%\label{tab:axioms}
\end{table}



In \cite{M93}, the semantics of if-then-else is given in terms of an action of $B$ on a $\lor$-semilattice with $\bot$, in the sense of Remark~\ref{rem:balgebra}.
Theorem~\ref{thm:new1} suggests that it might be interesting to investigate a semantics of if-then-else as a weakened (omitting (R5)) linear action of $B$ on an $\mathbb{F}_2$-vector space. 

\subsection{Boolean actions and sheaves}
In this section we give a sheaf representation of the structures of Theorem \ref{thm:new1}.

We start with some notation and preliminaries.

Given a Boolean algebra $B$ we ambiguously denote by $B$ the category whose objects are the elements of $B$ and there is an arrow $b \to c$ if $b \le c$. 
Let $\EuScript{V}$ be a variety of algebras. 
A sheaf $S: B^{\op} \to \EuScript{V}$ is a functor 
%(that is, a set $S(b)$ for every $b \in B$, and a function $S(b) \to S(c)$ when $c \le b$)
such that $S(b) \simeq S(c) \times S(d)$ when $b=c \lor d$ and $cd=0$. 
%The same applies when $S: B^{\op} \to \EuScript{V}$ is a sheaf with values in a variety of algebras. 

We recall Bergman's result concerning $B$-sets that will be used in the proof of Theorem \ref{thm:new2}. 

\begin{theorem} \cite[Theorem 12]{B91}
  \label{thm:berg}
  Let $B$ be a Boolean ring and let $\EuScript{V}$ be a variety of algebras.  
  \begin{itemize}
    \item For any set $X \neq \varnothing$, to equip $X$ with the structure of $B$-set amounts to give a sheaf $S:B^{\op} \to \mathbf{Set}$ such that $X \simeq S(1)$. 
    \item For any $\varnothing \neq X \in \EuScript{V}$, to equip $X$ with an action of $B$ amounts to give a sheaf $S:B^{\op} \to \EuScript{V}$ such that $X \simeq S(1)$. 
  \end{itemize}
\end{theorem}

Thanks to Theorem \ref{thm:new1} we obtain the main result of this section, which is a representation for models of an affine theory whose ring of coefficients is Boolean. 
\begin{theorem}
  \label{thm:new2}
    Let $B$ be a Boolean ring, and let $A_B$ its corresponding affine theory. 
    Let $X$ be a nonempty model of $A_B$ together with the choice of an element $o \in X$. 
    To equip $X$ with an action of $B$ amounts to give a sheaf $S:B^{\op} \to \mathbf{Vect}_{\mathbb{F}_2}$ such that $X \simeq S(1)$.
    %The category of pointed models of $A_B$, that is models $X$ of $A_B$ , is equivalent to the category of 
    %sheaves $S: B^{\op} \to \mathbf{Vect}_{\mathbb{F}_2}$ of  
    %$\mathbb{F}_2$-vector spaces over $B$ 
    %such that $X \simeq S(1)$. 
    %together with the choice of a point in the global sections of $S$, i.e., $o \in S(1)$. 
\end{theorem}
\begin{proof}
  Firstly, observe that by Theorem \ref{thm:berg} and by Remark \ref{rem:balgebra} a sheaf $S:B^{\op} \to \mathbf{Vect}_{\mathbb{F}_2}$ can be equivalently described as a $\mathbb{F}_2$-vector space $X$ such that: 
  \begin{itemize}
    \item $X$ is a $B$-set, i.e. there is an action of $B$ on $X$ satisfying (B1)-(B5); 
    \item $X$ satisfies (L1), i.e. $b(x+y,z+w) = b(x,z)+b(y,w)$ for every $b \in B$ and $x,y,z,w \in X$. 
  \end{itemize}
  Given a model $X$ of $A_B$ together with an element $o \in X$, we define an $\mathbb{F}_2$-vector space as in Theorem \ref{thm:new1}. 
  Now, $X$ has an action of $B$ that satisfies (R1)-(R4) and (L1) by Theorem \ref{thm:new1}.  
  By Lemma \ref{lem:rtob}, it is enough to prove that (R5) holds. 
  We prove, equivalently, that (B2) holds (here we prove the first part of (B2), the other part being similar).
  Using the usual algebraic manipulations involving the Mal'cev operation, we can see that:  
  \begin{equation}
    \label{eq:comb}
    b(x,y)=b(x,o)+(1-b)(y,o)\text{.}
  \end{equation}
  This implies that:
  % the action satisfies (R5) as well (here we prove part of (B2), the other part is similar): 
  \begin{align*}
    b(b(x,y),z) & = b(b(x,o)+(1-b)(y,o)) +(1-b)(z,o) & \eqref{eq:comb} \\
    & = b(b(x,o),o)+b((1-b)(y,o),o)+(1-b)(z,o) & \text{(L1)} \\
    & = b(x,o)+0(x,o)+(1-b)(z,o) & \text{(B5)} \\
    & = b(x,z) & \eqref{eq:comb}
  \end{align*} 
  as desiderd. 
  %We define an action of $B$ by letting $b(x,y):=bx+(1-b)y$. 
  %This action satisfies (B1)-(B5).
  The converse is obvious in light of Theorem \ref{thm:new1}. 
\end{proof}

\section{Conclusions and vistas}
\label{sec:conclusions}
This work has shown that affine and hyperaffine algebraic theories form a natural setting for representing commutative and Boolean rings, respectively.  
By developing a framework that accommodates both constructions, we have clarified the features that make affine theories the counterpart of commutative rings, and hyperaffine theories the counterpart of Boolean rings. 
Moreover, the study of the models of both reveals a close connection with the algebraic treatment of the if-then-else construct.
%, offering a new perspective on McCarthy’s axioms through the lens of algebraic theories.
Several directions arise.
A first step is to generalise the correspondence established here beyond rings, toward suitable classes of commutative \emph{semi}rings.
%Semirings are ubiquitous in the semantics of programming languages, particularly in quantitative, probabilistic, and resource-sensitive settings, and identifying the right algebraic theories corresponding to them could yield a unifying perspective on a wide spectrum of computational effects.
A second line of investigation concerns the algebras described in Theorem~\ref{thm:new1}, which suggest an equational foundation for an if-then-else operator in linear contexts, where the combination of programs obeys algebraic rather than set-theoretic structure.
%Developing this idea could bridge the gap between Boolean and linear semantics, leading to a more general theory of conditional computation.
Finally, a deeper understanding of the free algebras of affine theories
% could reveal new algebraic phenomena.
might lead %, inspired by \cite{BLPS18}, 
to a notion of $n$-dimensional commutative ring and 
to an extension of the classical theory of affine combinations into higher dimensions.
%although this remains speculative.
% it hints at rich interactions between algebraic theories, categorical geometry, and the semantics of computation—an interplay that we believe is only beginning to unfold.
%
% ---- Bibliography ----
%
% BibTeX users should specify bibliography style 'splncs04'.
% References will then be sorted and formatted in the correct style.
%
\bibliographystyle{abbrv}
\bibliography{csl}

\end{document}

Conversely, given a $\mathbb{F}_2$-vector space $X$ with an action of $B$ satisfying (B1)-(B5) and (L1),
  by Theorem \ref{thm:new1}, it remains to prove that 
  $(b+c)(x,y)=b(x,y)+c(x,y)$ for every $b,c \in B$ and $x,y \in X$. 
  Observe that, by (L2) and (B3), $a(x,y)=a(x+0,0+y)=a(x,0)+(1-a)(y,0)=a(x,0)+a(0,y)$;
  it is therefore enough to prove that $b(x,0)+c(x,0)=(b+c)(x,0)$ and $b(0,x)+c(0,x)=(b+c)(0,x)$ (we leave the latter to the reader).  
  %we extend the Abelian group structure $(X,+,0)$ to a $B$-module as follows. 
  Let $bx:=b(x,0)$ for every $b \in B$.  
  %the axioms of $B$-modules are satisfied: 
  First observe that 
  \begin{equation}
    \label{eq:smart}
    (1-c)(x,0)=c(0,x)=c(x+x,x+0)=c(x,x)+c(x,0)=x+cx
  \end{equation}
  by (B3), (L1) and (B1). 
  Therefore 
  \begin{align*}
    (b+c)x & = b(c(0,x),c(x,0)) & \text{(B3), (R4)}\\
    & = b(x+cx,0+cx) & \eqref{eq:smart} \\
    & =b(x,0) +b(cx,cx) & \text{(L1)} \\
    & =bx + cx & \text{(B1)} 
  \end{align*} 
  which gives the desired result. 


We start with some notation and preliminaries.

For a set $X$, we denote by $\Eq(X)$ the set of equivalence relations on $X$, by $\Delta_X:=\{(x,x) : x \in X\}$ the diagonal relation and by $\nabla_X:=X^2$ the total relation. 
An element $R \in \Eq(X)$ is called a \emph{factor relation} if there is $\lnot R \in \Eq(X)$ such that 
$R \cap \lnot R = \Delta_X$ and $R \circ \lnot R = \nabla_X \text{.}$
It is well-known (cf. \cite[Section 4.4]{ALV1}) that for a set $X$, there is a bijection between 
  %\begin{enumerate}
    (i) binary operations $a:X^2 \to X$ satisfying (R1) and (R5);
    %i.e. \emph{rectangular band} structures on $X$;  
    (ii) factor relations on $X$;
    (iii) the decompositions of $X$ into two factors: $X \simeq A \times B$. 

Given a Boolean algebra $B$ we ambiguously denote by $B$ the category whose objects are the elements of $B$ and there is an arrow $b \to c$ if $b \le c$. 
Then a sheaf $S: B^{\op} \to \mathbf{Set}$ is a functor 
%(that is, a set $S(b)$ for every $b \in B$, and a function $S(b) \to S(c)$ when $c \le b$)
such that $S(b) \simeq S(c) \times S(d)$ when $b=c \lor d$ and $cd=0$. 


We recall Bergman's result concerning $B$-sets that will be used in the proof of Theorem \ref{thm:new2}. 
%We take more space to contextualise and state Bergman's main result of \cite{B91}. 

  %\end{enumerate}
%The next result was established by Bergman in \cite{B91}.
%Bergman's next result can be seen as an enrichment of the foregoing fact. 

\begin{theorem} 
  [Lemma 10, Theorem 12, Lemma 16 \cite{B91}]
  %\label{thm:berg}
  Let $X$ be a set. 
  A set $F$ of binary operations on $X$ satisfying \emph{(R1), (R2)} and \emph{(R5)} gives rise to a Boolean ring $B$ and a $B$-set structure on $X$. 
  Moreover, a $B$-set structure on $X$ is equivalently given by 
  \begin{enumerate}
    \item a family of equivalence relations $\{R_b : b \in B\} \subseteq \Eq(X)$ satisfying: 
  \begin{enumerate}
    \item $R_b \circ R_{1-b} = \nabla_X$; 
    \item $b \le c \implies R_c \subseteq R_b$; 
    \item $R_b \cap R_c \subseteq R_{b \lor c}$; 
    \item $R_1 = \Delta_X$ 
  \end{enumerate}
  \item a sheaf $S: B^{\op} \to \mathbf{Set}$ of sets on $B$ such that $X \simeq S(1)$.  
\end{enumerate}
\end{theorem}

\begin{comment}
\begin{proof}
  Given a set $F$ of functions satisfying (R1),(R2) and (R5), if we define $(1-b)(x,y):=b(y,x)$ and $(bc)(x,y):=b(c(x,y),y)$ for every $b,c \in F$ we get a Boolean ring structure on $F$. 
  If $B$ is a Boolean ring acting on $X$, define $R_b:=\{(x,y) : b(x,y)=y\}$ to obtain the desired family of equivalence relations.  
  If $B$ is a Boolean ring, and $\{R_b : b \in B\}$ is a set of equivalence relations indexed by $B$ satisfying the properties (a)-(d), 
  construct $S: B^{\op} \to \mathbf{Set}$ defining $S(b):=X/R_b$. 
  %Recall that a basis for $\Stone(B)$ is given by $\{U_b : b \in B\}$ where $U_b:=\{U \in \Stone(B) : b \in U\}$. 
  %Then define $S(U_b):= X/\theta_b$.  
  Finally, let $S: B^{\op} \to \mathbf{Set}$ be a sheaf such that $X$ is the set of global sections of $S$.
  Then for each $b \in B$, the restriction maps $X \to S(b)$, $X \to S(1-b)$ constitute the projection maps of a product decomposition: $X \simeq S(b) \times S(1-b)$. 
  For each $b \in B$ we define the functions $b, 1-b: X^2 \to X$ to be the projections associated with this product 
  %(cf. Remark \ref{rem:band}). 
  to obtain a $B$-set structure on $X$. 
 
\end{proof}
\end{comment}

\begin{remark}
  \label{rem:balgebra2}
  More generally Bergman proves \cite[Theorem 12]{B91} that for any variety of algebras $\EuScript{V}$, there is an equivalence between the category of algebras with action of a Boolean ring $B$, and the category of sheaves $B^{\op} \to \EuScript{V}$.
\end{remark}

For a set $X$, we denote by $\Eq(X)$ the set of equivalence relations on $X$, by $\Delta_X:=\{(x,x) : x \in X\}$ the diagonal relation and by $\nabla_X:=X^2$ the total relation. 
An element $R \in \Eq(X)$ is called a \emph{factor relation} if there is $\lnot R \in \Eq(X)$ such that 
$R \cap \lnot R = \Delta_X$ and $R \circ \lnot R = \nabla_X \text{.}$
It is well-known (cf. \cite[Section 4.4]{ALV1}) that for a set $X$, there is a bijection between 
  %\begin{enumerate}
    (i) binary operations $a:X^2 \to X$ satisfying (R1) and (R5);
    %i.e. \emph{rectangular band} structures on $X$;  
    (ii) factor relations on $X$;
    (iii) the decompositions of $X$ into two factors: $X \simeq A \times B$. 

